\section*{Chapitre \ref{ch:m3:blockchains-for-traders} --- Blockchains pour les traders}

\begin{solution}{exo:m3:blockchains-for-traders:1}
$20 \times (1 + \tfrac18 \times 0.25/0.5) = 21.25$ gwei.
\end{solution}

\begin{solution}{exo:m3:blockchains-for-traders:2}
Environ une heure (six blocs de dix minutes en moyenne)~; environ 12.8 minutes (deux époques) pour la \omterm{def:m3:blockchains-for-traders:finality}{finalité} d'Ethereum.
\end{solution}

\begin{solution}{exo:m3:blockchains-for-traders:3}
L'ancrage est maintenu par l'arbitrage entre le rachat à un dollar et le cours de marché~; seuls les détenteurs
ayant un accès primaire peuvent réaliser cet arbitrage. Si peu d'acteurs peuvent racheter, ou si le rachat est fermé
(heures bancaires, une réserve gelée), le cours de marché peut tomber sous un dollar, sans rien pour le ramener.
\end{solution}

\begin{solution}{exo:m3:blockchains-for-traders:4}
Blocs pleins~: $\ln 2/\ln 1.125 = 5.9$, donc six blocs suffisent à plus que doubler les \omterm{def:m3:blockchains-for-traders:gas}{frais de base} ($1.125^6 = 2.03$). Blocs vides~:
$\ln 0.5/\ln 0.875 = 5.2$, donc cinq blocs ramènent le frais à 51~\% et six à 45~\%.
\end{solution}

\begin{solution}{exo:m3:blockchains-for-traders:5}
$150\,000 \times 32 \times 10^{-9} \times 3\,000 = \$14.40$.
\end{solution}

\begin{solution}{exo:m3:blockchains-for-traders:6}
Le \omterm{def:m3:blockchains-for-traders:rollup}{séquenceur} d'un \omterm{def:m3:blockchains-for-traders:rollup}{rollup} l'inclut tout de suite et à faible coût, mais la couche de base ne le garantit que lorsque le lot
est publié et, pour un \omterm{def:m3:blockchains-for-traders:rollup}{rollup} optimiste, que sa période de contestation est écoulée, ou, pour un \omterm{def:m3:blockchains-for-traders:rollup}{rollup} à validité, que sa preuve
est vérifiée~; jusque-là, l'ordre fixé par le \omterm{def:m3:blockchains-for-traders:rollup}{séquenceur} pourrait en principe être remis en cause.
\end{solution}

\begin{solution}{exo:m3:blockchains-for-traders:7}
Les \omterm{def:m3:blockchains-for-traders:gas}{frais de base} culminent à 39.20 gwei~; les blocs cessent d'être pleins au bloc 11, une fois que les \omterm{def:m3:blockchains-for-traders:gas}{frais de base} ont suffisamment
monté pour ramener la demande sous la limite.
\end{solution}

\begin{solution}{exo:m3:blockchains-for-traders:8}
L'ancrage suppose aussi que le rachat fonctionne quand les détenteurs le veulent~: heures bancaires, défaillance d'un dépositaire (comme à Silicon
Valley Bank en 2023), un gel juridique ou la solvabilité même de l'émetteur peuvent interrompre les rachats, et les cours secondaires
chutent alors, quelle que soit la valeur des réserves.
\end{solution}

\begin{solution}{pb:m3:blockchains-for-traders:1}
\textbf{1.} 5.40~\$. \textbf{2.} 11.25 gwei après un bloc, 105.45 gwei après vingt. \textbf{3.} 48.35~\$.
\textbf{4.} Environ quatre minutes (vingt créneaux de 12 secondes). \textbf{5.} Les \omterm{def:m3:blockchains-for-traders:gas}{frais de base} sont brûlés~; les \omterm{def:m3:blockchains-for-traders:gas}{frais de priorité} vont au
proposeur du bloc. \textbf{6.} $60/(150\,000 \times 10^{-9} \times 3\,000) - 2 = 131.33$ gwei. \textbf{7.} Après
22 blocs pleins ($\ln 13.13/\ln 1.125 = 21.9$). \textbf{8.} Le seuil de rentabilité monte à 264.67 gwei, atteint après 28
blocs pleins. \textbf{9.} Concurrence pour l'inclusion~: d'autres searchers enchérissent sur la même opportunité, et c'est le \omterm{def:m3:blockchains-for-traders:gas}{frais de priorité}, et non
le \omterm{def:m3:blockchains-for-traders:gas}{frais de base}, qui décide de l'ordre. \textbf{10.} Il baisse de 12.5~\% au plus par bloc, jusqu'à ce que la demande atteigne la cible.
\textbf{11.} Oui, mais ses exécutions ne sont définitives que lorsque les lots sont réglés, la liquidité et les opportunités diffèrent, et son
\omterm{def:m3:blockchains-for-traders:rollup}{séquenceur} contrôle l'ordre des transactions. \textbf{12.} Il masque la transaction aux autres searchers jusqu'à l'inclusion, au
prix de payer le constructeur de blocs plutôt que de rivaliser dans le \omterm{def:m3:blockchains-for-traders:mempool}{mempool} ouvert. \textbf{13.} Pour que les proposeurs ne puissent pas tirer profit
en le gonflant, et pour rendre les frais prévisibles~; leur brûlage profite aussi à tous les détenteurs. \textbf{14.} Leurs
transactions coûtent dix fois plus cher ou attendent. \textbf{15.} Quelles transactions et quels contrats sont en concurrence, et les \omterm{def:m3:blockchains-for-traders:gas}{frais de priorité}
versés par les autres. \textbf{16.} Il tarife la congestion de façon transparente et prévisible, mais il attribue l'espace aux plus
offrants, qui lors d'un pic sont ceux qui ont le plus à extraire. \textbf{17.} Combien de temps les transferts mettent pour devenir définitifs et
ce qu'ils coûtent en congestion, afin qu'une transaction entre plateformes soit évaluée avec les deux. \textbf{18.} Le \omterm{def:m3:blockchains-for-traders:gas}{gas} est un coût fixe par
transaction, donc il pèse davantage sur un petit profit. \textbf{19.} \emph{Résultat clé~:} vingt blocs pleins
multiplient le \omterm{def:m3:blockchains-for-traders:gas}{frais de base} par 10.5 (de 10 à 105.45 gwei)~; l'arbitrage de 60~\$ cesse d'être rentable après 22 blocs pleins.
\textbf{20.} Position et inclusion dans un bloc.
\end{solution}

\begin{solution}{iq:m3:blockchains-for-traders:1}
Le point à partir duquel une transaction ne peut plus être annulée. Une plateforme ne crédite les dépôts qu'après celui-ci, de sorte qu'un transfert
entre bourses est bloqué pendant ce temps et exposé aux prix~; un dépôt annulé est une perte pour la plateforme.
\iqlookfor{latence du règlement-livraison comme risque d'inventaire.}
\end{solution}

\begin{solution}{iq:m3:blockchains-for-traders:2}
Les \omterm{def:m3:blockchains-for-traders:gas}{frais de base} sont fixés par le protocole à partir du \omterm{def:m3:blockchains-for-traders:gas}{gas} consommé dans le bloc précédent (à la hausse ou à la baisse, de 12.5~\% au plus), puis ils sont brûlés~; les
\omterm{def:m3:blockchains-for-traders:gas}{frais de priorité} sont choisis par l'utilisateur et versé au proposeur pour être inclus avant les autres.
\iqlookfor{la règle de mise à jour, le brûlage et le \omterm{def:m3:blockchains-for-traders:gas}{frais de priorité} comme enchère.}
\end{solution}

\begin{solution}{iq:m3:blockchains-for-traders:3}
Arbitrage sur le marché primaire par des institutions qui émettent et rachètent à un dollar face aux cours secondaires. Il échoue si
les réserves perdent de leur valeur ou deviennent inaccessibles, si le rachat est suspendu ou restreint, ou si la solvabilité de l'émetteur ou
sa situation juridique est mise en doute.
\iqlookfor{le mécanisme et les conditions qu'il requiert.}
\end{solution}

\begin{solution}{iq:m3:blockchains-for-traders:4}
Le \omterm{def:m3:blockchains-for-traders:key}{jeton} bridgé est une créance sur des actifs bloqués dans le \omterm{def:m3:blockchains-for-traders:bridge}{bridge}~: bogues de \omterm{def:m3:blockchains-for-traders:contract}{smart contract}, clés de signature compromises, la
gouvernance du \omterm{def:m3:blockchains-for-traders:bridge}{bridge}, et \omterm{def:m3:blockchains-for-traders:stable}{décrochages} (depeg) entre le \omterm{def:m3:blockchains-for-traders:key}{jeton} bridgé et le \omterm{def:m3:blockchains-for-traders:key}{jeton} natif.
\iqlookfor{le \omterm{def:m3:blockchains-for-traders:bridge}{bridge} comme dépositaire concentré.}
\end{solution}

\begin{solution}{iq:m3:blockchains-for-traders:5}
Ne créditer qu'après la \omterm{def:m3:blockchains-for-traders:finality}{finalité} de la chaîne (ou un nombre de confirmations fixé selon le montant à risque et l'historique des réorganisations de la
chaîne), et afficher jusque-là la transaction comme en attente. Risques~: une réorganisation annule la transaction, une double dépense, ou une
bifurcation (fork).
\iqlookfor{politique de confirmation fondée sur le risque.}
\end{solution}

\begin{solution}{iq:m3:blockchains-for-traders:6}
Observer le \omterm{def:m3:blockchains-for-traders:mempool}{mempool} et les blocs récents~: la distribution des \omterm{def:m3:blockchains-for-traders:gas}{frais de priorité} des transactions incluses, bloc par bloc, la file en
attente et leurs \omterm{def:m3:blockchains-for-traders:gas}{frais de priorité}, et la trajectoire des \omterm{def:m3:blockchains-for-traders:gas}{frais de base}. Estimer, pour chaque niveau de \omterm{def:m3:blockchains-for-traders:gas}{frais de priorité}, la probabilité d'inclusion dans les deux
prochains blocs (\omterm{def:m3:blockchains-for-traders:gas}{frais de base} projeté, concurrence des transactions en attente et arrivées nouvelles attendues), et choisir les
\omterm{def:m3:blockchains-for-traders:gas}{frais de priorité} les plus bas permettant une probabilité supérieure à 95~\%~; mettre à jour à chaque bloc et procéder à un back-test sur les inclusions réalisées.
\iqlookfor{données, modèle probabiliste et calibration sur les résultats observés.}
\end{solution}
